\documentclass[10pt,a4paper]{amsart}
\usepackage[margin=19mm]{geometry}
\usepackage{amsmath,amssymb,mathtools}
\usepackage[scr=rsfs,cal=euler]{mathalfa}
\usepackage{xfrac}
\usepackage[unicode,hidelinks,bookmarks=false]{hyperref}
\newcommand{\ie}{i.e.\ }
\newcommand{\cf}{cf.\ }
\newcommand{\resp}{respectively\ }
\newcommand{\Subsection}{\subsection}
\providecommand{\nobreakdash}{\nobreak\hbox{-}}
\setlength{\emergencystretch}{2em}
\multlinegap0pt
\usepackage{amssymb}
\usepackage{float}
\newtheorem{theorem}{Theorem}
\newtheorem{lemma}{Lemma}
\newcommand{\ord}{\mathrm{ord}}
\title{Elements of finite order in the normalizer of a~maximal torus of a semisimple group}
\author{Ivan Arzhantsev, Alexey Galt,, Alexey Staroletov}
\date{Source: arXiv:2604.08108v1 (CC BY 4.0)}
\begin{document}
\maketitle

\noindent\emph{Source and licence.} Elements of finite order in the normalizer of a~maximal torus of a semisimple group. Version arXiv:2604.08108v1 (CC BY 4.0), distributed by arXiv under the Creative Commons Attribution 4.0 International licence, http://creativecommons.org/licenses/by/4.0/, as recorded in the arXiv licence metadata for this exact version and shown on the abstract page https://arxiv.org/abs/2604.08108v1. Full licence notice: \texttt{licenses/arxiv-2604.08108-CC-BY-4.0.txt}.\par\smallskip

\noindent\textit{Editorial scope.} Counted unit: the article's theorem tmain (source0.tex lines 206--218). The article's complete original proof of it is delivered with it (source0.tex lines 339--376, one \texttt{proof} environment of 38 lines, which the article places away from the statement). No mathematics is written, repaired or re-derived: the statement and the proof are the article's own lines, in the article's own notation, and no retained line is substituted or re-flowed.  Retained uncounted as the source-local context the proof needs: lines 327--329.  Nothing was omitted from any retained span, so the package carries no omission record.  No retained line contains a citation, so the wrapper carries no bibliography and no bibliography entry.  No retained line contains a figure environment, an image-inclusion command or drawing code, so no image is delivered and the package declares no asset dependency.  Editorial formatting only, in addition: an AMS \texttt{amsart} preamble that restates the article's own declarations and macros used by the retained lines, namely the environments \texttt{theorem}, \texttt{lemma} on separate counters, together with the standard packages those lines need.  The retained environments print in this document as Theorem 1, Lemma 1 in order of appearance, whereas the article prints the same items with the numbering of its own full document.  The complete native source of the article as distributed in the arXiv e-print for this exact version is bundled unchanged as \texttt{sources/198136/source0.tex}.
\begin{lemma} \label{lem}
The orbits of the action of the torus $T$ on the component $N_w$ by conjugation are subsets of the form $n_wK$, where $K$ is a coset of $T$ with respect to the subtorus $T(w)$. In particular, for a fixed $w\in W$, all these orbits are closed and have the same dimension.
\end{lemma}

\begin{theorem} \label{tmain}
Let $m(w)$ be the multiplicity of an element $w\in W$. Then the set $D_k(w)$ of elements of order $k$ in the component $N_w$ is:
\begin{itemize}
\item[a)]
empty if $d(w)$ does not divide $k$;
\item[b)]
a disjoint union of $m(w)\nu(s,\sigma(w))$ irreducible closed subvarieties $C_i$ if $d(w)=\ord(w)$ and $k=d(w)s$;
\item[c)]
a disjoint union of irreducible closed subvarieties $C_i$ if $d(w)=2\cdot\ord(w)$ and $k=d(w)s$, where the number of such subvarieties is
$m(w)(\nu(s,\sigma(w))+\nu(2s,\sigma(w)))$ for odd $s$ and $m(w)\nu(2s,\sigma(w))$ for even~$s$.
\end{itemize}
Each subvariety $C_i$ is isomorphic to a torus of dimension $r-\sigma(w)$ and is an orbit of the action of $T$ on $N_w$ by conjugation.
\end{theorem}

\begin{proof}[Proof of Theorem~\ref{tmain}]
Let $d=\ord(w)$ and assume that $n_w$ is an element of the smallest order in the component $N_w$. Recall that $d(w)$ denotes the order of $n_w$ and that $d(w)=d$ or $2d$.

Note that for any $t\in T$, we have
$$
(n_wt)^d=n_w^d\beta_w(t), \quad \text{where} \quad \beta_w(t)=(n_w^{1-d}tn_w^{d-1})\ldots(n_w^{-1}tn_w)t=t\varphi_w(t)\ldots\varphi_w^{d-1}(t).
$$
The map $\beta_w\colon T\to T$ is a homomorphism. Let us fix a basis in the character lattice of the torus $T$, and let $\Phi_w$, $A_w$, and $B_w$ be the matrices of the homomorphisms $\varphi_w$, $\alpha_w$, and $\beta_w$ in this basis, respectively. Since the addition of matrices corresponds to the multiplication of homomorphisms, we have
$$
A_w=\Phi_w^{-1}-E \quad \text{and} \quad B_w=E+\Phi_w+\ldots+\Phi_w^{d-1}.
$$
It is easy to see that $B_wA_w=0$. This means that the image $T(w)$ of $\alpha_w$ lies in the kernel of~$\beta_w$. Furthermore, the rank of the matrix $A_w$ is $r-\sigma(w)$. Since the matrix $\Phi_w$ has order~$d$, this matrix is diagonalizable and its eigenvalues $\epsilon$ are $d$-th roots of unity. The eigenvalues of the matrix $B_w$ are of the form
$$
1+\epsilon+\ldots+\epsilon^{d-1}.
$$
This number is equal to $d$ if $\epsilon=1$, and is $0$ otherwise. This shows that the rank of $B_w$ is~$\sigma(w)$. Thus, the subtorus $T(w)$ is the identity component of the kernel of $\beta_w$, and the image of this homomorphism is a subtorus $S(w)$ of dimension $\sigma(w)$ in $T$. By the definition of $m(w)$, the kernel of $\beta_w$ is the union of $m(w)$ cosets $K_i$ of the torus $T$ by the subtorus~$T(w)$.

\smallskip

Consider two cases. First, let $d(w)=d$. If $d$ does not divide $k$, then $(n_wt)^k$ does not lie in $T$. Consequently, this element cannot be equal to the identity, and the set $D_k(w)$ is empty. Next, we assume $k=ds$. In this case, $(n_wt)^k=\beta_w(t)^s$. Thus, the element $n_wt$ lies in $D_k(w)$ if and only if $\beta_w(t)$ is an element of order $s$ in the subtorus $S(w)$. There are exactly $\nu(s,\sigma(w))$ such elements. Therefore, the set $D_k(w)$ is identified with the union of $m(w)\nu(s,\sigma(w))$ subvarieties $C_i=n_wK_i$, where $K_i$ are cosets of $T$ with respect to the subtorus $T(w)$. According to Lemma~\ref{lem}, these subvarieties coincide with the orbits of the action of $T$ on $N_w$ by conjugation.

\smallskip

Now let $d(w)=2d$. Then $t_0:=n_w^d$ is an element of order two in $T$. This element does not lie in the subtorus $S(w)$. Indeed, if $t_0=\beta_w(t)$, then
$$
(n_wt)^d=n_w^d\beta_w(t)=t_0t_0=e.
$$
This contradicts the condition $d(w)=2d$.

If $d$ does not divide $k$, then $(n_wt)^k$ does not lie in $T$ and the set $D_k(w)$ is empty. Next, we assume $k=dl$. Then $(n_wt)^k=t_0^l\beta_w(t)^l$. If $l$ is odd, the element $(n_wt)^k=t_0\beta_w(t)^l$ cannot be equal to the identity, and the set $D_k(w)$ is again empty. Therefore, we further assume that $k=d(w)s$, so $(n_wt)^k=\beta_w(t)^{2s}$.
It follows that $\beta_w(t)^{2s}=1$. If, moreover, the order of $n_wt$ is strictly less than $k$, then by what has been proven, this order is equal to $2ds'$ for some proper divisor $s'$ of $s$. In this case, $\beta_w(t)^{2s'}=1$. Consequently, if the order of $n_wt$ is $k$, the element $\beta(t)$ either has order $2s$ in $S(w)$ or has order $s''$, where the divisor $s''$ of $2s$ does not divide any proper even divisor of $2s$. The latter is possible only if $s$ is odd and $s''=s$.

Thus, in the case of even $s$, we find that the set $D_k(w)$ is the union of $m(w)\nu(2s,\sigma(w))$ subsets $C_i=n_wK_i$, where $K_i$ are cosets of $T$ with respect to $T(w)$. For odd $s$, the number of subsets $C_i=n_wK_i$ is ${m(w)(\nu(s,\sigma(w))+\nu(2s,\sigma(w)))}$: the first summand in this sum corresponds to the case $(n_wt)^{ds}=t_0$, and the second to the case $(n_wt)^{ds}=t_0s_0$, where $s_0$ is an element of order two in $S(w)$.

\smallskip

Theorem~\ref{tmain} is proved.
\end{proof}

\end{document}
